\chapter{Physics Before Formulas}

Today's global financial market is arguably one of the largest engineered systems in existence: millions of instruments, fragmented venues, market makers, brokers, clearing systems, 
data centers, scripts, and autonomous agents interacting in nanoseconds. 
The New York Stock Exchange, which started when 24 brokers signed the Buttonwood Agreement in 1792, has grown into a global network of exchanges and trading venues that operate around the clock. 
In the modern era, the conception of a "market" has expanded to include not only traditional stock exchanges but also electronic trading platforms, dark pools, decentralized finance (DeFi) networks, prediction markets.

When viewing the market as an independent patchwork of venues, participants, and instruments, 
it is easy to get lost in the details. The goal of this chapter is the opposite: 
to abstract the market into a system that can be reasoned about, and then use that 
abstraction to guide the design of trading strategies that are robust by construction, 
not by accident.

This chapter provides a broad overview to see the forest through the trees: 
terminology first, then market edge (participants and workflows), market core (venues, execution, and data flow), 
and finally the intelligence layer (models, agents, risk controls, and failure modes).

\section{What Is the Market?}

When most people hear ``the market,'' they imagine a building, a bell, or a chart on a screen. 
That picture is not wrong, but it is incomplete enough to hurt you if you trade from it. 
The market is a network of interacting agents and venues that collectively determine prices and allocate capital, each participant operating under their own objectives, constraints, and information.

To start, we will define the market from a high level perspective, focusing on its structure, participants, and the flow of information and capital.
The second half of the chapter will focus on the market from a system engineering perspective, where we explore the market as a complex adaptive system that can be modeled, analyzed, and optimized.

\subsection{The Market's Structure}
The New York Stock Exchange traces back to the Buttonwood Agreement, signed by 24 brokers in New York on May 17, 1792.
In modern terms, Buttonwood functioned like a protocol: it defined who could trade with whom and standardized commissions. 
The point was trust and control, not speed. By requiring a broker for every trade, the agreement ensured that every transaction passed through a verified member who was personally responsible for the deal's success.

Prior to the Buttonwood Agreement, trading was far more informal, with brokers negotiating deals directly with each other in coffee houses and other ad hoc venues.
The lack of structure often led to disputes, inconsistent pricing, and a fear of \textbf{counterparty risk}: there was no guarantee that the person on the other side of a trade would actually deliver the stock or the cash.

By formalizing the rules, the Buttonwood signers helped establish trust and confidence in the financial system. In many ways, the agreement functioned as an early whitelist, where brokers served as transparent validator nodes in a distributed network.
That newfound confidence proved crucial for growth: it enabled more efficient price discovery, increased liquidity, and attracted broader participation from investors.
The market expanded in lockstep with the American economy as more companies sought to raise capital and more investors looked for opportunity.

In the early days, investors did not directly communicate with the exchange. Instead, they relied on brokerage offices or written instructions to place orders.
On the exchange floor, brokers gathered around designated trading posts shouting bids and offers. Only when two brokers agreed on a price was the trade recorded in notebooks.

Information about prices and trades then had to be transmitted back to investors through messengers, couriers, or telegraph updates. This process introduced unavoidable delays between activity on the exchange floor and the information available to distant investors.

In effect, the market had encountered its first \textbf{latency arbitrage} problem: participants with faster access to information could exploit those with slower access.
An investor located in Boston, for example, might have to wait hours or even days for a courier to deliver trade information from New York, while a trader on the floor could act on that information immediately.
The manual communication infrastructure that supported early trading had been designed for a much smaller market, and it struggled to keep pace with the accelerating volume of transactions.
This latency gap meant an investor might attempt to buy a stock at a certain price, but by the time their order reached the exchange, the price could already have changed, resulting in a missed opportunity or a less favorable trade.

To address this gap, Edward A. Calahan invented the ticker tape machine, which transmitted stock prices over existing telegraph infrastructure.
The device was one of the earliest examples of \textit{near-real-time data dissemination}.
Ticker tape machines functioned as specialized telegraphs whose sole purpose was to print stock symbols, prices, and volume onto a narrow paper strip. 
By 1929, ticker tape machines were installed in brokerage offices, investor's homes, and financial institutions across the United States, providing a continuous stream of market data to investors and traders.
The result was a dramatic improvement in the speed and accuracy of information flow from Wall Street to the rest of the country. For the first time, investors could make decisions based on near-real-time data, rather than relying on delayed reports.
In many ways, the ticker tape transformed financial markets from a localized trading floor into a distributed information network.
Ticker tape was not without its physical limitations. Machines could only print a limited number of characters per second, and the paper strips were prone to jamming and tearing.
Under normal conditions these issues were manageable, and by 1929 improved models could relay data at 285 characters per minute. But the system had a hard ceiling: when trading volume outpaced print speed, the tape fell behind, and every investor reading it was making decisions on stale prices.

From a systems perspective, the stock market had become an information network whose demand for data transmission was beginning to exceed its available bandwidth.
October 1929 proved the point catastrophically. As trading volume surged past what the ticker infrastructure could handle, reporting delays grew from minutes to hours, and a feedback loop of panic selling and stale data amplified the crash far beyond what the underlying economic conditions alone would have caused.

\begin{casehistory}[The 1929 Crash as an Information-System Failure]
Leading up to October 1929, several macroeconomic and political factors were already placing pressure on the financial system. In August 1929, the Federal Reserve raised interest rates from 5\% to 6\%, tightening credit conditions across the economy.
Financial forecaster Roger Babson had publicly warned of a coming collapse, an event that later became known as the ``Babson Break.''
At the same time, international instability, including the Hatry Crisis in the United Kingdom, caused British investors to withdraw capital from American markets.
Domestically, income inequality meant that a relatively small portion of the population controlled most investment capital, limiting the pool of new buyers needed to sustain rising prices.
Debates surrounding the proposed Smoot--Hawley Tariff Act also created fears of a global trade conflict.

These factors increased market volatility and trading activity, pushing the ticker tape system toward its operational limits.

During the week of the crash, trading volume surged to unprecedented levels. On October 21, 1929, approximately six million shares were traded, overwhelming ticker systems and causing reporting delays of nearly 100 minutes.
By October 24, later known as \textit{Black Thursday}, trading volume reached 12.9 million shares, and ticker delays expanded to nearly four hours.

For investors relying on ticker tape updates, the information they were receiving was no longer current market data, but rather a delayed record of trades that had occurred hours earlier.
Investors observing falling prices on the ticker believed the market was collapsing in real time, even though the prices they were seeing reflected conditions from much earlier in the trading session.
As panic selling accelerated, the volume of trades increased even further, widening the gap between actual market activity and the data being disseminated to the public.

By October 28 (\textit{Black Monday}), trading volume reached 9.2 million shares with ticker delays approaching three hours.
The following day, October 29, 1929, would become known as \textit{Black Tuesday}: 16.4 million shares were traded, an unprecedented level of activity for the era.

The infrastructure responsible for distributing price data could no longer keep pace with the rate at which new trades were occurring.
What began as an economic correction quickly escalated into a systemic information failure, where millions of investors were making decisions based on data that was already hours out of date.
\end{casehistory}

Viewed through a modern lens, the events of October 1929 illustrate a critical principle: price discovery is fundamentally constrained by the speed at which information can propagate through the system.
The crash did not just expose a limitation of ticker tape; it revealed that any market whose data infrastructure cannot keep pace with its trading activity is structurally fragile.
Solving this problem would require more than faster printers. It would require a fundamentally different approach to how prices were generated, matched, and distributed.

\subsection{The Digital Revolution and the NASDAQ Era}

The answer, when it finally arrived, was electronic execution. NASDAQ launched in 1971 as the first fully electronic stock market, replacing the physical ticker with digital price quotation and, eventually, automated order matching.

Where ticker tape had been a one-way broadcast---prices printed onto paper and shipped outward---electronic systems closed the loop. Quotes, orders, and confirmations could travel the same network in both directions, collapsing the gap between observing a price and acting on it.
Over time, electronic routing, matching, and reporting became normal across venues. The same information bottleneck that had amplified the 1929 crash was, in principle, eliminated: digital systems could scale throughput in ways that mechanical printers never could.

Digital systems also made market data more machine-readable and more available at scale, which is exactly what quantitative research needed. But computerization did not magically ``clean'' the data. It introduced new forms of noise, latency effects, and cross-venue inconsistency that quants still have to model explicitly.

\subsection{The Market as a Real-Time Decision Field}

In this course, we treat the market as a high-dimensional, real-time decision environment.

Every tick is a decision event. It may be generated by a human, an automated strategy, or a hybrid workflow. Either way, each print is an action taken under constraints, not a timeless truth.

\subsection{Primary vs Secondary Markets}

Primary markets are where new securities are issued (for example, IPO allocation and initial sale). Secondary markets are where existing claims are continuously repriced and transferred.

Most daily equity trading you see is secondary-market activity. That flow is less about raising new capital for firms and more about ongoing price discovery among holders with competing views.

\subsection{Core Market Participants}

The high-level player map from lecture is the right starting point:

\begin{itemize}
	\item Dealers/market makers: provide liquidity, quote two-sided prices, and manage inventory risk.
	\item Brokers: route and execute customer orders, with explicit obligations around execution quality.
	\item Speculators (including retail, quants, and students): seek returns from directional, statistical, or structural edges.
\end{itemize}

If you only look at exchanges and ignore these actors, you are not modeling the market. You are modeling a dashboard.

\subsection{Dealers and Market Making}

Dealers do not behave like long-only investors. Their core business is facilitating flow while managing inventory and spread risk.

In practice, many realized fills are the product of a dealer's willingness to hold risk briefly in exchange for spread economics and flow scale.

\subsection{Brokers and Order Routing}

Brokers are the transport layer between client intent and executable flow. Route choice, venue selection, and timing directly affect realized outcomes.

"Best execution" is not a slogan. It is an operational constraint with measurable consequences for slippage, fill quality, and strategy PnL.

\subsection{Financial Intermediaries}

Beyond brokers and dealers, large intermediaries (banks, funds, and asset managers) pool capital, transform exposures, and repackage risk.

These institutions change both the shape of order flow and the risk topology of the system. Ignoring them leads to naive causal stories.


\section{System Lens}

\sidenote{Anchor each implementation step to market mechanics first.}
